% Standalone problem document, generated from the adjacent README.md.
% Compile with XeLaTeX. Mathematical content is maintained in Markdown.
\documentclass[11pt,a4paper]{article}
\usepackage[fontsize=10.5pt]{scrextend}
\usepackage[margin=22mm,headheight=15pt,headsep=9mm,footskip=11mm]{geometry}
\usepackage{fontspec}
\setmainfont{texgyrepagella}[Extension=.otf,UprightFont=*-regular,BoldFont=*-bold,ItalicFont=*-italic,BoldItalicFont=*-bolditalic]
\setsansfont{texgyreheros}[Extension=.otf,UprightFont=*-regular,BoldFont=*-bold,ItalicFont=*-italic,BoldItalicFont=*-bolditalic]
\setmonofont{texgyrecursor}[Extension=.otf,UprightFont=*-regular,BoldFont=*-bold,ItalicFont=*-italic,BoldItalicFont=*-bolditalic,Scale=MatchLowercase]
\usepackage{amsmath,amssymb}
\usepackage{unicode-math}
\setmathfont{texgyrepagella-math.otf}
\usepackage{xcolor,microtype,enumitem,needspace,fancyhdr}
\usepackage{bookmark}
\definecolor{ink}{HTML}{17344B}
\definecolor{accent}{HTML}{197D80}
\definecolor{muted}{HTML}{596773}
\hypersetup{colorlinks=true,linkcolor=accent,urlcolor=accent,pdftitle={SP-05 - Symmetric
minimizer for a positive definite Jordan--Kronecker
product},pdfauthor={Open Problems in Numerical Linear Algebra}}
\urlstyle{same}
\setlength{\parindent}{0pt}
\setlength{\parskip}{5pt plus 1pt minus 1pt}
\linespread{1.04}
\setlength{\emergencystretch}{3em}
\widowpenalty=10000
\clubpenalty=10000
\displaywidowpenalty=10000
\allowdisplaybreaks[1]
\setlist{nosep,leftmargin=1.5em}
\providecommand{\tightlist}{\setlength{\itemsep}{0pt}\setlength{\parskip}{0pt}}
\providecommand{\pandocbounded}[1]{#1}
\pagestyle{fancy}
\fancyhf{}
\fancyhead[L]{\sffamily\footnotesize\color{muted}OPEN PROBLEMS IN NUMERICAL LINEAR ALGEBRA}
\fancyhead[R]{\sffamily\bfseries\footnotesize\color{accent}SP-05}
\fancyfoot[L]{\parbox[b]{0.9\textwidth}{\sffamily\fontsize{8}{10}\selectfont\color{muted}Editorial ratings; a bounded search cannot certify that a problem remains unsolved.\\Literature check: 2026-09-22}}
\fancyfoot[R]{\sffamily\footnotesize\color{muted}\thepage}
\renewcommand{\headrulewidth}{0.3pt}
\renewcommand{\headrule}{\hbox to\headwidth{\color{accent}\leaders\hrule height \headrulewidth\hfill}}
\makeatletter
\renewcommand{\section}{\Needspace{5\baselineskip}\@startsection{section}{1}{0pt}{12pt plus 2pt minus 2pt}{4pt}{\sffamily\bfseries\large\color{ink}}}
\renewcommand{\subsection}{\Needspace{4\baselineskip}\@startsection{subsection}{2}{0pt}{9pt plus 1pt minus 1pt}{3pt}{\sffamily\bfseries\normalsize\color{ink}}}
\makeatother
\setcounter{secnumdepth}{0}
\begin{document}
{\sffamily\small\color{muted}Eigenvalues and inverse problems\par}
\vspace{3pt}
{\raggedright\hyphenpenalty=10000\exhyphenpenalty=10000\sffamily\bfseries\fontsize{21}{24}\selectfont\color{ink}Symmetric
minimizer for a positive definite Jordan--Kronecker product\par}
\vspace{7pt}
{\sffamily\small\textbf{Difficulty:} challenging\quad\textbf{Importance:} interesting
to the community\par}
{\sffamily\small\bfseries\color{accent}Status: Lean verified\par}
\vspace{4pt}
{\color{accent}\hrule height 0.8pt}
\vspace{6pt}
\textbf{Topic:} Structured eigenvalue problems and semidefinite
optimization.

\textbf{Rating rationale:} Challenging reflects a structured
minimum-eigenvalue comparison for arbitrary positive definite pairs;
community impact connects Kronecker eigenproblems and semidefinite
optimization.

\subsection{Resolution --- 2026-09-11}\label{resolution-2026-09-11}

\textbf{Affirmative resolution by Matthew J. Colbrook} (Department of
Applied Mathematics and Theoretical Physics, University of Cambridge).
See the
\href{https://github.com/ajt60gaibb/OpenProblemsInNLA/blob/main/eigenvalues-and-inverse-problems/SP-05/solution.md}{complete
manuscript}, \textbf{Theorem SP-05, sections 1--3}
(\href{https://github.com/ajt60gaibb/OpenProblemsInNLA/blob/main/eigenvalues-and-inverse-problems/SP-05/solution.pdf}{PDF}
·
\href{https://github.com/ajt60gaibb/OpenProblemsInNLA/blob/main/eigenvalues-and-inverse-problems/SP-05/solution.tex}{LaTeX}),
prepared 11 September 2026.

For arbitrary real symmetric positive definite \(A,B\), a nonzero real
positive-semidefinite eigenmatrix attains the smallest eigenvalue of
\(X\mapsto AXB+BXA\). Section 3 derives the exact
symmetric/skew-symmetric Rayleigh-quotient inequality in the original
target, without commutativity, rank restrictions or a simple-eigenvalue
assumption.

The original proof draft was generated in a ChatGPT conversation. A
separate Codex agent independently verified the full proof and its match
to the exact target on 11 September 2026:
\href{https://github.com/ajt60gaibb/OpenProblemsInNLA/blob/main/references/colbrook-2026-09-11/verification/reviews/SP-05-review.md}{detailed
PASS review}. The review records a hash of the unchanged proof text.
This is independent agent verification, not external human peer review
or formal certification.
\href{https://github.com/ajt60gaibb/OpenProblemsInNLA/blob/main/references/colbrook-2026-09-11/README.md}{The
submission history and diagnostic record} preserve the initial solution
claim. The ratings above are historical, and the earlier literature
checks below are retained.

\subsection{Lean verification --- 22 September
2026}\label{lean-verification-22-september-2026}

The full original inequality is formally proved for every dimension at
least two and every real symmetric positive-definite pair. The proof
constructs a nonzero real PSD global minimizing eigenmatrix and
establishes that both actual sector minima exist and are attained. All
four exports passed LeanCert kernel audits, two independent nonauthor
final reviews, and
\href{https://github.com/sidneyholden1/OpenProblemsInNLA/actions/runs/35693379965}{fresh
isolated Linux Comparator/default-kernel verification} at the
\href{https://github.com/sidneyholden1/OpenProblemsInNLA/tree/8596dc32d33c931de352aace53df8f02a67602ce/eigenvalues-and-inverse-problems/SP-05/lean}{immutable
proof revision}.

Mathematical resolution: Matthew J. Colbrook. Original conjecture:
Kalantarova and Tunçel. Formalization: Sidney Holden with OpenAI Codex
assistance. The retained checks include rejection and sandbox controls;
AI-agent reviews are not external human peer review.
\href{https://github.com/sidneyholden1/OpenProblemsInNLA/blob/main/eigenvalues-and-inverse-problems/SP-05/lean/README.md}{Proof
and reviews} ·
\href{https://github.com/sidneyholden1/OpenProblemsInNLA/blob/main/docs/lean/verification-2026-09-22/SP-05/README.md}{Linux
evidence}.

\subsection{Problem statement}\label{problem-statement}

Let \(n\ge 2\) and let \(A,B\in\mathbb R^{n\times n}\) be symmetric
positive definite. Let \(T\in\mathbb R^{n^2\times n^2}\) be the
commutation matrix defined by
\(T\mathop{\mathrm{vec}}\nolimits(X)=\mathop{\mathrm{vec}}\nolimits(X^T)\),
where \(\mathop{\mathrm{vec}}\nolimits\) stacks columns. Is it always
true that

\[\min_{\substack{u\in\mathbb R^{n^2}\setminus\{0\}\\Tu=u}}
\frac{u^T(A\otimes B)u}{u^Tu}
\;\le\;
\min_{\substack{w\in\mathbb R^{n^2}\setminus\{0\}\\Tw=-w}}
\frac{w^T(A\otimes B)w}{w^Tw}?\]

Equivalently, must the smallest eigenvalue of \(A\otimes B+B\otimes A\)
have an eigenvector \(\mathop{\mathrm{vec}}\nolimits(X)\) with
\(X=X^T\ne0\)?

\subsection{Why it matters}\label{why-it-matters}

An affirmative answer would identify which invariant subspace contains
the smallest eigenvalue of an important structured matrix operator,
relevant to positivity checks in semidefinite optimization.

\subsection{References}\label{references}

\begin{itemize}
\tightlist
\item
  N. Kalantarova and L. Tunçel,
  \href{https://arxiv.org/abs/1805.09737}{On the spectral structure of
  Jordan--Kronecker products of symmetric and skew-symmetric matrices},
  \emph{Linear Algebra and its Applications} 608 (2021), 343--362.
  Conjecture 1, equation (7), p.~12 of arXiv v3; the preceding
  discussion explains the optimization connection.
\item
  D. Kressner and B. Vandereycken,
  \href{https://arxiv.org/abs/2608.20875}{A counterexample to the
  symmetric-maximizer conjecture for Lyapunov operators}, 2026,
  abstract. This resolves a different symmetric-maximizer conjecture and
  is included to distinguish the recent result from this problem.
\end{itemize}

\subsection{Earlier status check ---
2026-09-08}\label{earlier-status-check-2026-09-08}

On 2026-09-08, checked arXiv v3 (2020-08-08, still the latest listed
version) and searched combinations of ``Jordan-Kronecker'', ``positive
definite'', ``conjecture'', ``smallest'', ``counterexample'', and the
authors' names, including 2025--2026 results. No resolution of
Conjecture 1 was located. The source's counterexamples to general
symmetric-matrix interlacing do not dispose of its separately stated
positive definite conjecture. The August 2026 Lyapunov result concerns
an operator norm maximizer, rather than this positive definite
Jordan--Kronecker minimum. No recent explicit reaffirmation of this
exact conjecture was found; the status evidence is therefore bounded.

\subsection{Audit update --- 2026-09-10}\label{audit-update-2026-09-10}

Rechecked the
\href{https://www.math.uwaterloo.ca/~ltuncel/publications/1805.09737.pdf}{author
manuscript}, Conjecture 1 and equation (7). Its positive definite,
minimum-eigenvalue formulation matches this entry. Searches for the
Jordan--Kronecker symmetric-minimizer conjecture and later
counterexamples found no resolution. Results for indefinite inputs or a
maximum-eigenvalue objective do not settle the displayed claim.
\end{document}
